% ============================================================================
% CHAPITRE 3 : INTRODUCTION
% ============================================================================

% ----------------------------------------------------------------------------
\section{INTRODUCTION}
\label{sec:introduction_constructions}
% ----------------------------------------------------------------------------

\begin{sloppypar}
Une borne d'impossibilit\'e ne dit rien, \`a elle seule, de ce qu'il est
possible de construire. Le chapitre pr\'ec\'edent a \'etabli que le produit
$\Conf\cdot\Rel\cdot\Opp$ ne peut d\'epasser une quantit\'e gouvern\'ee par
la capacit\'e d'extraction institutionnelle. Deux questions restent
ouvertes. La premi\`ere porte sur le caract\`ere atteignable de cette borne~; la
seconde demande si la preuve qui l'\'etablit r\'esiste \`a
une v\'erification m\'ecanique.

Ce chapitre r\'ealise les actes \textbf{S} et \textbf{C} du cycle
canonique. La section~\ref{sec:sotp_base} construit le Semantic One-Time
Pad, prototype qui montre que le paradigme Gold-Iron-Clay est
instanciable. La section~\ref{sec:sotp_construction} g\'en\'eralise ce
prototype en une famille de trois protocoles, chacun taill\'e pour un
r\'egime d'entropie contextuelle, puis formule la conjecture de leur
optimalit\'e universelle. La section~\ref{sec:lean4} soumet le noyau alg\'ebrique du
trilemme \`a l'assistant de preuve Lean4. La
section~\ref{sec:conclusion_constructions} dresse le bilan de ces deux
actes. L'acte \textbf{X}, exp\'erimenter, fait l'objet du
chapitre~\ref{chap:experimentation}.

% ----------------------------------------------------------------------------
\subsection{D\'emarche m\'ethodologique~: les actes \textbf{S} et \textbf{C}}
\label{subsec:demarche_saturer_certifier}
% ----------------------------------------------------------------------------

\paragraph{Saturer (\textbf{S}).}
Une borne d'impossibilit\'e prend toute sa valeur op\'erationnelle
lorsqu'on sait l'approcher. La m\'ethodologie vis\'ee est celle de la
construction explicite, qui associe \`a chaque r\'egime d'entropie
contextuelle un protocole de la famille SOTP (SOTP-Low, SOTP-Base, SOTP-High)
propos\'e comme construction de r\'ef\'erence. Leurs propri\'et\'es
cryptographiques sont prouv\'ees~; la saturation de la borne, qui porte sur
la fiabilit\'e de reconstruction, reste une conjecture
(Remarque~\ref{rem:statut_sotp}).

\paragraph{Certifier (\textbf{C}).}
Un th\'eor\`eme qui croise trois disciplines peut receler des hypoth\`eses
implicites que l'intuition math\'ematique ne d\'etecte pas. La
m\'ethodologie Lean4 impose une explicitation compl\`ete, puisque chaque
hypoth\`ese non prouv\'ee dans le syst\`eme doit \^etre \'enonc\'ee
comme propri\'et\'e ou comme pr\'emisse. Cette exigence a conduit \`a
expliciter l'hypoth\`ese de r\'egime et \`a pr\'eciser trois \'enonc\'es
publi\'es.

\begin{remark}[Statut des r\'esultats de ce chapitre]\label{rem:statut_sotp}
Les constructions S-OTP et SOTP sont \'evalu\'ees sur l'\'echelle
op\'erationnelle CLO, qui combine confidentialit\'e statistique, fiabilit\'e de
v\'erification $\Relver$ (Remarque~\ref{rem:rel_ver}) et opposabilit\'e
sous des conventions d'interpr\'etation explicites. Ces m\'etriques ne sont
pas compar\'ees \`a la borne du Th\'eor\`eme~\ref{thm:trilemme_cro_conditionnel},
qui porte sur $\Relrec$ et \'enonce que, pour tout protocole \`a confidentialit\'e
statistique $1-\negl(\lambda)$, $\Relrec \leq (1+\kappa/\lambda)/\lambda$.
L'existence de protocoles qui approchent la borne raffin\'ee~\eqref{eq:trilemme_cro_raffine}
est formul\'ee comme conjecture (Conjecture~\ref{conj:optimalite_universelle}).
Les composantes de preuve qui ne d\'ependent que du contexte public $C$
n'apportent aucune information conditionnelle, puisque
$I(V; g(C) \mid C) = 0$ pour toute fonction $g$. La couche Clay est
donc analys\'ee comme un \emph{format de pr\'esentation} qui facilite
l'interpr\'etation (baisse du co\^ut cognitif) sans ajouter
d'information au sens de $\Opp$. Les expressions num\'eriques de
saturation (facteurs $0{,}75$, $0{,}5$, $0{,}45$) sont des heuristiques de
conception, sans statut de minoration prouv\'ee.
\end{remark}
\end{sloppypar}
